
We develop posterior sampling at the level of the general interpolant path of Section~\ref{subsec:flow_models}, so one construction covers SI, FM, and diffusion (DM-SDE) and yields both SDE and ODE samplers. Assume the prior model trained, supplying a velocity $\fmvel_\tau$, score $\genscore_\tau$, or SI drift for the path onto $\conddist{\bx_1}{\bx_0}$; we want samples from the one-step posterior $\conddist{\bx_1}{\obs,\bx_0}$ of Eq.~\eqref{eq:posterior_distribution} without retraining. Conditioning on $\obs$ preserves the path structure and only retargets it, so a single guidance term yields the whole family of posterior samplers (SI-SDE, DM-SDE, guided ODE), which share their observation handling and differ by a scalar weight.
\subsection{Posterior sampling: SDE and ODE}
\label{subsec:posterior_dynamics}

We want an SDE (or ODE) whose terminal law is the one-step posterior $\conddist{\bx_1}{\obs,\bx_0}$ of Eq.~\eqref{eq:posterior_distribution}, built from a prior model that is already trained and never retrained. Conditioning on $\obs$ preserves the interpolant path of Section~\ref{subsec:flow_models} and only reweights it by the \emph{likelihood tilt}
\begin{align}\label{eq:likelihood_tilt}
    \Phi_\tau^{\obs}(\bx,\bx_0) := \conddist{\obs}{\bX_\tau=\bx,\bx_0}
    = \E\!\left[\conddist{\obs}{\bx_1}\,\middle|\,\bX_\tau=\bx,\bX_0=\bx_0\right],
\end{align}
the observation probability given the current state. This shifts the prior score and velocity by the likelihood score $\nabla_{\bx}\log\Phi_\tau^{\obs}$ (Appendix~\ref{appendix:proof_conditioning}), which is all that conditioning changes and which a single guidance term carries through to every sampler below.

\begin{theorem}[Unified posterior sampler]\label{theorem:unified_posterior}
Let the prior model define an interpolant path with velocity $\fmvel_\tau$ and score $\genscore_\tau$ whose terminal law is $\conddist{\bx_1}{\bx_0}$, and let $\Phi_\tau^{\obs}$ be the likelihood tilt \eqref{eq:likelihood_tilt}. For any diffusion schedule $\gdiff_\tau\ge0$, the SDE
\begin{align}\label{eq:unified_posterior_drift}
    \rd\bXstar_\tau
    = \underbrace{\left[\fmvel_\tau(\bXstar_\tau,\bx_0) + \tfrac12\gdiff_\tau^2\,\genscore_\tau(\bXstar_\tau,\bx_0)\right]}_{\text{prior drift }\drift^{\gdiff}_\tau}\rd\tau
    + \underbrace{\big(\vscoef_\tau + \tfrac12\gdiff_\tau^2\big)}_{=:\,\gweight_\tau}\nabla_{\bXstar_\tau}\log\Phi_\tau^{\obs}(\bXstar_\tau,\bx_0)\,\rd\tau
    + \gdiff_\tau\,\rd\bW_\tau ,
\end{align}
initialised from $\bXstar_0\sim p_0^{\obs}$ with $p_0^{\obs}(\rd\bx)\propto\Phi_0^{\obs}(\bx,\bx_0)\,p_0(\rd\bx)$, has marginals $p_\tau^{\obs}(\bx\mid\bx_0)\propto\Phi_\tau^{\obs}\,p_\tau$ for all $\tau\in[0,1]$, and in particular generates $\bXstar_1\sim\conddist{\bx_1}{\obs,\bx_0}$.
\end{theorem}

The dynamics split into the prior drift and a single likelihood-score correction with guidance weight $\gweight_\tau=\vscoef_\tau+\tfrac12\gdiff_\tau^2$ (the velocity--score term plus the diffusion term, $\vscoef_\tau$ from Eq.~\eqref{eq:general_score_velocity}). In the ideal setting (exact prior and tilt) every member of the family yields exact posterior samples; the proof and the full conditioning argument are in Appendix~\ref{appendix:proof_posterior}. The prior drift, the likelihood score, and the observation handling are shared across samplers, which differ only through $\gdiff_\tau$ and hence $\gweight_\tau$:
\begin{align}\label{eq:three_drifts}
\begin{array}{l@{\;}l@{\quad}l}
\text{SI-SDE} &(\gdiff_\tau=\gamma_\tau): &
  \drift^{\gamma}_\tau\,\rd\tau
  +\ \big(\vscoef_\tau+\tfrac12\gamma_\tau^2\big)\,\nabla\log\Phi_\tau^{\obs}\,\rd\tau\ +\ \gamma_\tau\,\rd\bW_\tau, \\[4pt]
\text{DM-SDE} &(\gdiff_\tau>0): &
  \drift^{\gdiff}_\tau\,\rd\tau
  +\ \big(\vscoef_\tau+\tfrac12\gdiff_\tau^2\big)\,\nabla\log\Phi_\tau^{\obs}\,\rd\tau\ +\ \gdiff_\tau\,\rd\bW_\tau, \\[4pt]
\text{FM-ODE} &(\gdiff_\tau=0): &
  \fmvel_\tau\,\rd\tau
  +\ \vscoef_\tau\,\nabla\log\Phi_\tau^{\obs}\,\rd\tau .
\end{array}
\end{align}
At $\gdiff_\tau=0$ the noise and the prior-drift score term vanish, leaving the deterministic guided ODE; $\gdiff_\tau>0$ restores both, so SI, denoising diffusion (DM-SDE), and the guided flow are one family. An alternative posterior SDE from Doob's $h$-transform conditions trajectories rather than marginals and has no deterministic member; we use the marginal construction \eqref{eq:unified_posterior_drift} throughout (Appendix~\ref{appendix:proof_posterior}).






% \newpage
% We present the posterior sampling method at the level of the general affine Gaussian path of Section~\ref{subsec:flow_models}, so the same construction applies to SI, FM, and DM (referred to as DM-SDE) and yields both SDE and ODE samplers. The prior model is assumed trained, providing a velocity $\fmvel_\tau$ (equivalently a score $\genscore_\tau$) or drift (in the SI case) for the path that transports a tractable source to the forecast prior $\conddist{\bx_1}{\bx_0}$; the aim is to sample from the one-step posterior $\conddist{\bx_1}{\obs,\bx_0}$ of Eq.~\eqref{eq:posterior_distribution} without retraining. The central observation is that conditioning on $\obs$ leaves the affine Gaussian structure intact and only retargets it, so a \emph{single} guidance term produces a family of posterior samplers -- SI-SDE, DM-SDE, and the guided ODE -- with identical observation handling, differing only in a scalar weight.

% \subsection{Posterior sampling: SDE and ODE}
% \label{subsec:posterior_dynamics}

% \paragraph{Conditioning the affine path.}
% Let $\obs = \obsoperator(\bx_1) + \obsnoise$ with likelihood $\conddist{\obs}{\bx_1}$, and define the \emph{likelihood tilt}
% \begin{align}\label{eq:likelihood_tilt}
%     \Phi_\tau^{\obs}(\bx,\bx_0) := \conddist{\obs}{\bX_\tau=\bx,\bx_0}
%     = \E\!\left[\conddist{\obs}{\bx_1}\,\middle|\,\bX_\tau=\bx,\bX_0=\bx_0\right],
% \end{align}
% the probability of the observation given the current state, with expectation over the path coupling $\conddist{\bx_1}{\bx_\tau,\bx_0}$. By Bayes' rule the posterior marginals factor as
% \begin{align}\label{eq:posterior_marginal}
%     p_\tau^{\obs}(\bx\mid\bx_0) := \conddist{\bx_\tau=\bx}{\obs,\bx_0}
%     = \frac{\Phi_\tau^{\obs}(\bx,\bx_0)\,p_\tau(\bx\mid\bx_0)}{\conddist{\obs}{\bx_0}},
% \end{align}
% so the posterior path is the prior path reweighted by the likelihood tilt, again an affine Gaussian path with the \emph{same} schedules and target replaced by the posterior. All proofs of this section are in Appendices~\ref{appendix:proof_conditioning}--\ref{appendix:proof_interpolant_likelihood}.

% \begin{proposition}[Posterior path]\label{prop:conditional_path}
% Conditioning the interpolant \eqref{eq:general_interpolant} on $\obs$ yields the affine Gaussian path
% \begin{align}
%     \bx_\tau = \alpha_\tau\bm a_0 + \sigpath_\tau\latentnoise + \beta_\tau\bx_1,
%     \qquad \bx_1\sim\conddist{\bx_1}{\obs,\bx_0},\ \ \latentnoise\sim\mathcal{N}(0,\bI),
% \end{align}
% with $\latentnoise$ still standard normal and independent of $\bx_1$. Consequently its score and velocity are the prior score and velocity shifted by the likelihood score,
% \begin{align}
%     \genscore_\tau^{\obs}(\bx,\bx_0)
%     &= \genscore_\tau(\bx,\bx_0) + \nabla_{\bx}\log\Phi_\tau^{\obs}(\bx,\bx_0),
%        \label{eq:conditional_score}\\
%     \fmvel_\tau^{\obs}(\bx,\bx_0)
%     &= \fmvel_\tau(\bx,\bx_0) + \vscoef_\tau\,\nabla_{\bx}\log\Phi_\tau^{\obs}(\bx,\bx_0),
%        \label{eq:conditional_velocity}
% \end{align}
% where $\vscoef_\tau$ is the velocity--score coefficient \eqref{eq:vscoef_def}.
% \end{proposition}

% The score relation \eqref{eq:conditional_score} is the standard Bayesian decomposition of score-based conditional sampling \cite{song_score-based_2021}; the velocity relation \eqref{eq:conditional_velocity} follows by applying the velocity--score duality \eqref{eq:general_velocity_of_score} to the posterior path, whose affine part (schedules and anchor) is unchanged by conditioning and cancels in $\fmvel_\tau^{\obs}-\fmvel_\tau$, leaving $\vscoef_\tau\nabla_{\bx}\log\Phi_\tau^{\obs}$.

% \paragraph{The unified posterior sampler.}
% The posterior path has its own marginal-preserving SDE family, from Proposition~\ref{prop:marginal_preservation} applied to $p_\tau^{\obs}$.

% \begin{theorem}[Unified posterior sampler]\label{theorem:unified_posterior}
% Let the prior model define an affine Gaussian path with velocity $\fmvel_\tau$ and score $\genscore_\tau$ whose terminal law is $\conddist{\bx_1}{\bx_0}$, and let $\Phi_\tau^{\obs}$ be the likelihood tilt \eqref{eq:likelihood_tilt}. For any diffusion schedule $\gdiff_\tau\ge0$, the SDE
% \begin{align}\label{eq:unified_posterior_sde}
%     \rd\bXstar_\tau
%     = \left[\fmvel_\tau^{\obs}(\bXstar_\tau,\bx_0) + \tfrac12\gdiff_\tau^2\,\genscore_\tau^{\obs}(\bXstar_\tau,\bx_0)\right]\rd\tau
%     + \gdiff_\tau\,\rd\bW_\tau,
%     \qquad \bXstar_0\sim p_0^{\obs},
% \end{align}
% initialised from the reweighted source $p_0^{\obs}(\rd\bx)\propto\Phi_0^{\obs}(\bx,\bx_0)\,p_0(\rd\bx)$, has marginals $p_\tau^{\obs}$ for all $\tau\in[0,1]$, and in particular generates terminal states $\bXstar_1\sim\conddist{\bx_1}{\obs,\bx_0}$.
% \end{theorem}

% This is the exactness statement: in the ideal setting (exact prior velocity/score and likelihood tilt) every member of the family produces \emph{exact} posterior samples (Appendix~\ref{appendix:proof_posterior}). Substituting \eqref{eq:conditional_score}--\eqref{eq:conditional_velocity} into \eqref{eq:unified_posterior_sde} separates the dynamics into the prior drift and a single likelihood-score correction,
% \begin{align}\label{eq:unified_posterior_drift}
%     \rd\bXstar_\tau
%     = \underbrace{\left[\fmvel_\tau(\bXstar_\tau,\bx_0) + \tfrac12\gdiff_\tau^2\,\genscore_\tau(\bXstar_\tau,\bx_0)\right]}_{\text{prior drift }\drift^{\gdiff}_\tau}\rd\tau
%     + \underbrace{\left(\vscoef_\tau + \tfrac12\gdiff_\tau^2\right)}_{=:\,\gweight_\tau}\nabla_{\bXstar_\tau}\log\Phi_\tau^{\obs}(\bXstar_\tau,\bx_0)\,\rd\tau
%     + \gdiff_\tau\,\rd\bW_\tau .
% \end{align}
% The correction pushes towards high observation likelihood while the prior drift pulls towards the prior, together steering the marginals from $p_0^{\obs}$ to the posterior. The single \emph{guidance weight}
% \begin{align}\label{eq:guidance_weight}
%     \gweight_\tau = \vscoef_\tau + \tfrac12\gdiff_\tau^2
% \end{align}
% is the only quantity that changes between samplers: $\vscoef_\tau$ from the velocity--score duality, $\tfrac12\gdiff_\tau^2$ from the chosen diffusion.

% \paragraph{SI, FM, and ODE: identical form, one locus of difference.}
% The SI-SDE, lifted DM-SDE, and guided FM-ODE are literally Eq.~\eqref{eq:unified_posterior_drift} at different $\gdiff_\tau$. Writing $\bar S_\tau:=\nabla_{\bXstar_\tau}\log\Phi_\tau^{\obs}$ for the shared likelihood score, the three drifts align as
% \begin{align}\label{eq:three_drifts}
% \begin{array}{l@{\;}l@{\quad}c@{\quad}l}
% \text{SI-SDE} &(\gdiff_\tau=\gamma_\tau): &
%   \drift^{\gamma}_\tau\,\rd\tau
%   &+\ \big(\vscoef_\tau+\tfrac12\gamma_\tau^2\big)\,\bar S_\tau\,\rd\tau\ +\ \gamma_\tau\,\rd\bW_\tau, \\[4pt]
% \text{DM-SDE} &(\gdiff_\tau>0): &
%   \drift^{\gdiff}_\tau\,\rd\tau
%   &+\ \big(\vscoef_\tau+\tfrac12\gdiff_\tau^2\big)\,\bar S_\tau\,\rd\tau\ +\ \gdiff_\tau\,\rd\bW_\tau, \\[4pt]
% \text{FM-ODE} &(\gdiff_\tau=0): &
%   \fmvel_\tau\,\rd\tau
%   &+\ \phantom{\big(}\vscoef_\tau\phantom{+\tfrac12\gdiff_\tau^2\big)}\,\bar S_\tau\,\rd\tau.
% \end{array}
% \end{align}
% The prior drift, the likelihood score $\bar S_\tau$, and the observation handling are common; the differences are exactly the diffusion $\gdiff_\tau$, the weight $\gweight_\tau$, and (through $p_0^{\obs}$) the source. Setting $\gdiff_\tau=0$ removes the noise and prior-drift score term, leaving the deterministic guided ODE; $\gdiff_\tau>0$ adds them back -- SI, denoising-diffusion (DM-SDE), and the guided flow are one family.

% \begin{remark}[Relation to Doob's $h$-transform]\label{rem:doob}
% An alternative posterior \emph{SDE} fixes $\gdiff_\tau>0$ and applies Doob's $h$-transform with $h(\bx)=\conddist{\obs}{\bx}$, giving the correction $\gdiff_\tau^2\nabla\log\Phi_\tau^{\obs}$ (Theorem~\ref{theorem:posterior_sde}, Appendix~\ref{appendix:proof_posterior}). This matches the posterior terminal law but conditions trajectories rather than marginals, so its intermediate marginals are not $p_\tau^{\obs}$ and it has no deterministic member. The marginal construction \eqref{eq:unified_posterior_drift} instead reproduces $p_\tau^{\obs}$ at every $\tau$ and carries the weight $\gweight_\tau$ consistent with the tractable interpolated-observation likelihood of Section~\ref{subsec:obs_interpolation}; we adopt it as the working sampler.
% \end{remark}

\subsection{Approximating the likelihood score}
\label{subsec:obs_interpolation}

The correction in \eqref{eq:unified_posterior_drift} needs the likelihood score $\nabla_{\bx}\log\Phi_\tau^{\obs}$ at intermediate states, but the exact tilt
\begin{align}\label{eq:phi_integral}
    \Phi_{\tau}^{\obs}(\bx_\tau, \bx_0)
    = \int \conddist{\obs}{\bx_1}\, \conddist{\bx_1}{\bx_\tau, \bx_0}\, \rd \bx_1
\end{align}
is intractable: it requires forward integration of the dynamics and backpropagation through the trajectory. We instead interpolate the observation, mirroring the state interpolant in observation space,
\begin{align} \label{eq:interpolated_observation}
    \interpolantobs_\tau(\bx_0, \obs) = \alpha_\tau \obsmatrix\bx_0 + \beta_\tau \obs,
    \qquad \interpolantobs_0 = \obsmatrix\bx_0,\quad \interpolantobs_1 = \obs,
\end{align}
which reduces to $\beta_\tau\obs$ for FM ($\alpha_\tau=0$). Replacing the tilt with the surrogate $\bar{\Phi}_{\tau}^{\interpolantobs_\tau}(\bx_\tau, \bx_0) = \conddist[\tau]{\interpolantobs_\tau}{\bx_\tau, \bx_0}$ gives the working form of \eqref{eq:unified_posterior_drift} for all three samplers,
\begin{align} \label{eq:interpolant_posterior_drift}
    \postdrift_\theta(\bx_{\tau}, \bx_0, \interpolantobs_\tau, \tau)
    =
    \drift^{\gdiff}_\tau(\bx_{\tau}, \bx_0)
    + \gweight_\tau\,
    \nabla_{\bx_\tau} \log\bar{\Phi}_{\tau}^{\interpolantobs_\tau}(\bx_\tau, \bx_0),
\end{align}
with interpolant-likelihood score $\bar s_\tau := \nabla_{\bx_\tau}\log\bar\Phi_\tau^{\interpolantobs_\tau}$. To evaluate the interpolant-likelihood score, we need the mean and covariance:
% Writing the zero-mean source as $\src_\tau := \sigpath_\tau\latentnoise$, so $\bx_\tau = \alpha_\tau\bx_0 + \beta_\tau\bx_1 + \src_\tau$, the surrogate has a closed form for both model classes.
\begin{lemma}[Interpolated observation likelihood]\label{lem:interpolated_likelihood}
Let $\obsoperator$ be linear with matrix $\obsmatrix$ and $\obsnoise \sim \mathcal{N}(0,\obscov)$. Then
\begin{equation}\label{eq:ybar_decomp}
  \interpolantobs_\tau = \obsmatrix\bx_\tau + \interpolantobsnoise_\tau,
  \qquad
  \interpolantobsnoise_\tau := \beta_\tau\obsnoise - \obsmatrix\src_\tau,
\end{equation}
where $\src_\tau$ is noise associated with the interpolant and $\conddist{\interpolantobs_\tau}{\bx_0, \bx_\tau}$ has mean and covariance
\begin{align}
  \bar{\mu}_\tau(\bx_\tau, \bx_0)
  = \obsmatrix\bx_\tau - \obsmatrix\,\srcmean_\tau(\bx_\tau,\bx_0),
     \quad % \\[3pt]
  \bar{\Sigma}_{\tau}(\bx_\tau, \bx_0)
  = \beta_\tau^2\obscov + \obsmatrix\,\srccov_\tau(\bx_\tau,\bx_0)\,\obsmatrix^\top,
\end{align}
with source moments $\srcmean_\tau := \E[\src_\tau\mid\bx_\tau,\bx_0]$ and $\srccov_\tau := \cov(\src_\tau\mid\bx_\tau,\bx_0)$.
\end{lemma}
The proof is in Appendix~\ref{appendix:proof_interpolant_likelihood}. The source moments are computable from the trained model, the two classes differing only in $\src_\tau$.



% \newpage

% \subsection{Approximating the likelihood score}
% \label{subsec:obs_interpolation}

% Evaluating the correction in \eqref{eq:unified_posterior_drift} requires the likelihood score $\nabla_{\bx}\log\Phi_\tau^{\obs}$ at intermediate states. The exact tilt,
% \begin{align}\label{eq:phi_integral}
%     \Phi_{\tau}^{\obs}(\bx_\tau, \bx_0)
%     = \int \conddist{\obs}{\bx_1}\, \conddist{\bx_1}{\bx_\tau, \bx_0}\, \rd \bx_1,
% \end{align}
% is intractable, requiring forward integration of the dynamics and backpropagation through the trajectory. We avoid this by interpolating the observation itself. Define the \emph{observation interpolant}
% \begin{align} \label{eq:interpolated_observation}
%     \interpolantobs_\tau(\bx_0, \obs) = \alpha_\tau \obsmatrix\bm{a}_0 + \beta_\tau \obs,
%     \qquad \interpolantobs_0 = \obsmatrix\bm a_0,\quad \interpolantobs_1 = \obs,
% \end{align}
% which mirrors the state interpolant in observation space, transporting $\obsmatrix\bm a_0$ at $\tau=0$ to $\obs$ at $\tau=1$ (for SI, $\interpolantobs_\tau = \alpha_\tau\obsmatrix\bx_0 + \beta_\tau\obs$; for FM, $\interpolantobs_\tau = \beta_\tau\obs$). The tilt is then replaced by the computable surrogate $\bar{\Phi}_{\tau}^{\interpolantobs_\tau}(\bx_\tau, \bx_0) = \conddist[\tau]{\interpolantobs_\tau}{\bx_\tau, \bx_0}$, giving the posterior drift correction in the interpolant-likelihood score $\bar S_\tau := \nabla_{\bx_\tau}\log\bar\Phi_\tau^{\interpolantobs_\tau}$,
% \begin{align} \label{eq:interpolant_posterior_drift}
%     \postdrift_\theta(\bx_{\tau}, \bx_0, \interpolantobs_\tau, \tau)
%     =
%     \drift^{\gdiff}_\tau(\bx_{\tau}, \bx_0)
%     + \gweight_\tau\,
%     \nabla_{\bx_\tau} \log\bar{\Phi}_{\tau}^{\interpolantobs_\tau}(\bx_\tau, \bx_0).
% \end{align}
% This is the working form of \eqref{eq:unified_posterior_drift} for all three samplers. Writing the zero-mean source process as $\src_\tau := \sigpath_\tau\latentnoise$ (so $\bx_\tau = \alpha_\tau\bm a_0 + \beta_\tau\bx_1 + \src_\tau$), we need the conditional law of $\interpolantobs_\tau$ given $\bx_\tau$, which holds for both model classes.

% \begin{lemma}[Interpolated observation likelihood]\label{lem:interpolated_likelihood}
% Let $\obsoperator$ be linear with matrix $\obsmatrix$ and the observation noise be Gaussian, $\obsnoise \sim \mathcal{N}(0,\obscov)$. Then the interpolated observation satisfies
% \begin{equation}\label{eq:ybar_decomp}
%   \interpolantobs_\tau = \obsmatrix\bx_\tau + \interpolantobsnoise_\tau,
%   \qquad
%   \interpolantobsnoise_\tau := \beta_\tau\obsnoise - \obsmatrix\src_\tau,
% \end{equation}
% and the conditional mean and covariance of $\conddist{\interpolantobs_\tau}{\bx_0, \bx_\tau}$ are
% \begin{align}
%   \E[\interpolantobs_\tau \mid \bx_\tau, \bx_0]
%   &= \obsmatrix\bx_\tau - \obsmatrix\,\srcmean_\tau(\bx_\tau,\bx_0) =: \bar{\mu}_\tau(\bx_\tau, \bx_0),
%      \label{eq:bias} \\[3pt]
%   \cov(\interpolantobs_\tau \mid \bx_\tau, \bx_0)
%   &= \beta_\tau^2\obscov + \obsmatrix\,\srccov_\tau(\bx_\tau,\bx_0)\,\obsmatrix^\top =: \bar{\Sigma}_{\tau}(\bx_\tau, \bx_0),
%      \label{eq:cov_correction}
% \end{align}
% where $\srcmean_\tau := \E[\src_\tau\mid\bx_\tau,\bx_0]$ and $\srccov_\tau := \cov(\src_\tau\mid\bx_\tau,\bx_0)$ are the conditional moments of the source process.
% \end{lemma}
% The proof is in Appendix~\ref{appendix:proof_interpolant_likelihood}. The source moments are computable from the trained model, the two classes differing only in $\src_\tau$.



\begin{lemma}[Source conditional moments]\label{lem:source_moments}
    The conditional moments of the source process are:
    \begin{itemize}
        \item \emph{Stochastic interpolants}, $\src_\tau = \gamma_\tau\bW_\tau$:
        \begin{align}
            \srcmean_\tau = -\gamma_\tau^2\tau A_\tau\!\left(\beta_\tau\drift_\theta(\bx,\bx_0,\tau) - c_\tau(\bx,\bx_0)\right),
            \quad
            \srccov_\tau = \gamma_\tau^2\tau\bI + \gamma_\tau^4\tau^2 A_\tau\!\left(\beta_\tau J_{\drift_\theta} - \dot\beta_\tau\bI\right),
        \end{align}
        with $A_\tau, c_\tau$ as in Eq.~\eqref{eq:SI_score} and $J_{\drift_\theta}=\nabla_{\bx}\drift_\theta$.
        \item \emph{Flow matching and diffusion model}, $\src_\tau = \alpha_\tau\latent$:
        \begin{align}\label{eq:fm_source_moments}
            \srcmean_\tau = -\alpha_\tau^2\,\genscore_\tau(\bx),
            \qquad
            \srccov_\tau = \alpha_\tau^2\bI + \alpha_\tau^4\,\nabla_{\bx}\genscore_\tau(\bx),
        \end{align}
        with $\genscore_\tau$ in closed form from Eq.~\eqref{eq:fm_score}.
    \end{itemize}
    \end{lemma}
    Both follow from Stein's and the second-order Tweedie identities (Appendix~\ref{appendix:proof_interpolant_likelihood}). In each case the conditional mean equals the observation interpolant at the model denoiser, $\bar\mu_\tau = \alpha_\tau\obsmatrix\bx_0 + \beta_\tau\obsmatrix\,\E[\bx_1\mid\bx_\tau,\bx_0]$, so Eq.~\eqref{eq:interpolant_posterior_drift} is a Tweedie-style guidance term shared by all three samplers.
    
    Taking the interpolant likelihood Gaussian,
    \begin{align}\label{eq:interpolant_likelihood_gaussian}
        \tilde{p}(\interpolantobs_\tau \mid \bx_\tau, \bx_0)
        := \mathcal{N}\!\left(\bar{\mu}_\tau(\bx_\tau, \bx_0),\, \bar{\Sigma}_\tau(\bx_\tau, \bx_0)\right),
    \end{align}
    the interpolant-likelihood score is, in closed form,
    \begin{align}\label{eq:interpolant_score_closed}
        \bar{S}_\tau
        = (\nabla_{\bx_\tau}\bar{\mu}_\tau)^\top\,\bar{\Sigma}_\tau^{-1}\,(\interpolantobs_\tau - \bar{\mu}_\tau),
    \end{align}
    with the covariance $\bar{\Sigma}_\tau$ held fixed in $\bx_\tau$ when differentiating, as is common in guidance methods. The accuracy turns on $\bar S_\tau$ using the inflated covariance $\bar{\Sigma}_\tau = \beta_\tau^2\obscov + \obsmatrix\srccov_\tau\obsmatrix^\top$ of Eq.~\eqref{eq:cov_correction}, with the full source covariance $\srccov_\tau$ of Lemma~\ref{lem:source_moments}. This is the $\Pi$GDM-style inflation \cite{song_pseudoinverse_2023} of the raw observation covariance $\beta_\tau^2\obscov$, available at no extra cost since the trained model already supplies $\srccov_\tau$; the inflation carries the prior-covariance information that lets the guidance reach the unobserved degrees of freedom. Diffusion Posterior Sampling (DPS) \cite{chung_diffusion_2023, chung_improving_2024, song_score-based_2021} drops it, using $\beta_\tau^2\obscov$ alone, which our experiments show loses substantial accuracy on sparse observations (Section~\ref{sec:results}).



% \begin{lemma}[Source conditional moments]\label{lem:source_moments}
% The conditional moments of the source process are:
% \begin{itemize}
%     \item \emph{Stochastic interpolants}, $\src_\tau = \gamma_\tau\bW_\tau$:
%     \begin{align}
%         \srcmean_\tau = -\gamma_\tau^2\tau A_\tau\!\left(\beta_\tau\drift_\theta(\bx,\bx_0,\tau) - c_\tau(\bx,\bx_0)\right),
%         \quad
%         \srccov_\tau = \gamma_\tau^2\tau\bI + \gamma_\tau^4\tau^2 A_\tau\!\left(\beta_\tau J_{\drift_\theta} - \dot\beta_\tau\bI\right),
%     \end{align}
%     with $A_\tau, c_\tau$ as in Eq.~\eqref{eq:SI_score} and $J_{\drift_\theta}=\nabla_{\bx}\drift_\theta$.
%     \item \emph{Flow matching and diffusion model}, $\src_\tau = \alpha_\tau\latent$:
%     \begin{align}\label{eq:fm_source_moments}
%         \srcmean_\tau = -\alpha_\tau^2\,\genscore_\tau(\bx),
%         \qquad
%         \srccov_\tau = \alpha_\tau^2\bI + \alpha_\tau^4\,\nabla_{\bx}\genscore_\tau(\bx),
%     \end{align}
%     with $\genscore_\tau$ given in closed form by Eq.~\eqref{eq:fm_score}.
% \end{itemize}
% \end{lemma}
% Both follow from Stein's and the second-order Tweedie identities (Appendix~\ref{appendix:proof_interpolant_likelihood}). In each case the conditional mean equals the observation interpolant at the model denoiser, $\bar\mu_\tau = \alpha_\tau\obsmatrix\bm a_0 + \beta_\tau\obsmatrix\,\E[\bx_1\mid\bx_\tau,\bx_0]$, so Eq.~\eqref{eq:interpolant_posterior_drift} is a Tweedie-style guidance term shared by all three samplers (SI via Wiener-source moments, FM via Tweedie moments through $\genscore_\tau$).

% Assuming the interpolant likelihood is Gaussian,
% \begin{align}\label{eq:interpolant_likelihood_gaussian}
%     \tilde{p}(\interpolantobs_\tau \mid \bx_\tau, \bx_0)
%     := \mathcal{N}\!\left(\bar{\mu}_\tau(\bx_\tau, \bx_0),\, \bar{\Sigma}_\tau(\bx_\tau, \bx_0)\right),
% \end{align}
% the interpolant-likelihood score is then, in closed form,
% \begin{align}\label{eq:interpolant_score_closed}
%     \bar{S}_\tau := \nabla_{\bx_\tau}\log\bar{\Phi}_\tau^{\interpolantobs_\tau}
%     = (\nabla_{\bx_\tau}\bar{\mu}_\tau)^\top\,\bar{\Sigma}_\tau^{-1}\,(\interpolantobs_\tau - \bar{\mu}_\tau),
% \end{align}
% where the gradients of the Gaussian log-density hold the covariance $\bar{\Sigma}_\tau$ fixed in $\bx_\tau$, as is common in guidance methods. The crucial feature is that $\bar S_\tau$ uses the \emph{inflated} covariance
% \begin{align}\label{eq:inflated_covariance}
%     \bar{\Sigma}_\tau = \beta_\tau^2\,\obscov + \obsmatrix\,\srccov_\tau\,\obsmatrix^\top
% \end{align}
% of Eq.~\eqref{eq:cov_correction}, with the \emph{full} source covariance $\srccov_\tau$ of Lemma~\ref{lem:source_moments}. This is exactly the $\Pi$GDM-style inflation \cite{song_pseudoinverse_2023} of the raw observation covariance $\beta_\tau^2\obscov$ -- here available for free, since $\srccov_\tau$ is already supplied by the trained model -- and is the locus of the method's accuracy: the inflation injects the prior-covariance information that lets the guidance reach the unobserved degrees of freedom. The cruder Diffusion Posterior Sampling (DPS) surrogate \cite{chung_diffusion_2023, chung_improving_2024, song_score-based_2021} drops the inflation, using the uninflated covariance $\beta_\tau^2\obscov$ alone; our experiments show this loses substantial accuracy on sparse observations (Section~\ref{sec:results}).

\subsection{Covariance approximation.} 
\label{subsec:covariance_approximation}

The inflated $\bar\Sigma_\tau$ carries a network-Jacobian term through $\srccov_\tau$. Dropping it replaces $\srccov_\tau$ by its isotropic part $\rho_\tau\bI$ ($\rho_\tau=\gamma_\tau^2\tau$ for SI, $\alpha_\tau^2$ for FM), giving the state-independent, precomputable
\begin{align}\label{eq:cov_cheap}
    \bar{\Sigma}_\tau \approx \beta_\tau^2\,\obscov + \rho_\tau\,\obsmatrix\obsmatrix^\top,
\end{align}
accurate when $\sigpath_\tau^2\|\nabla_{\bx}^2\log p_\tau\|\ll1$ (near $\tau\to1$, or near-Gaussian priors) and the method of choice for dense observations and large $N_y$. Retaining the full $\srccov_\tau$ needs $N_y$ Jacobian--vector products per ensemble member per step, $\mathcal{O}(E\,N_y)$ and prohibitive at field scale; the ensemble-shared approximation evaluates the source-covariance Jacobian once at the ensemble mean, factorising a single $N_y\times N_y$ system to $\mathcal{O}(N_y)$, exact as the ensemble tightens. This EnKF-analogue step recovers accuracy in the sparse regime where the Jacobian-free covariance degrades (Section~\ref{sec:results}).



\subsection{Comparison of posterior samplers}
\label{subsec:summary_table}\label{subsec:supplying_diffusion}

The family \eqref{eq:interpolant_posterior_drift} contains three samplers of practical interest, sharing the observation-interpolant likelihood with inflated covariance (Section~\ref{subsec:obs_interpolation}) and differing only in $\gdiff_\tau$, the weight $\gweight_\tau=\vscoef_\tau+\tfrac12\gdiff_\tau^2$, and the source. The SI-SDE ($\gdiff_\tau=\gamma_\tau$) reuses the native SDE \eqref{eq:SI_SDE}, whose drift $\drift_\theta=\fmvel_\tau+\tfrac12\gamma_\tau^2\genscore_\tau$ already contains the score, from the point-mass source $\delta_{\bx_0}$. The DM-SDE ($\gdiff_\tau>0$) lifts the FM ODE by recovering the score via \eqref{eq:fm_score}, initialised exactly from $\mathcal{N}(0,\bI)$ since $\Phi_0^{\obs}$ is constant (Appendix~\ref{appendix:proof_conditioning}). The guided ODE ($\gdiff_\tau=0$) is the probability-flow ODE of the posterior path with weight $\vscoef_\tau$, coinciding with interpolant-guided flow for linear inverse problems \cite{yan_fig_2024, pokle_training-free_2024}; being deterministic it is well posed for FM but not for SI, whose point-mass source maps to a single point. Table~\ref{tab:samplers} contrasts the three; the model-specific bookkeeping (source process, how velocity and score are read from the network, the resulting moments and $\bar\Sigma_\tau$) is in Table~\ref{tab:si_vs_fm} (Appendix~\ref{appendix:bookkeeping}).

\begin{table}[ht]
    \centering
    \caption{The three posterior samplers as members of the unified family \eqref{eq:unified_posterior_drift}. All share the prior drift $\drift^{\gdiff}_\tau=\fmvel_\tau+\tfrac12\gdiff_\tau^2\genscore_\tau$, the likelihood-score correction $\gweight_\tau\,\bar S_\tau$, and the observation-interpolant machinery; they differ only in $\gdiff_\tau$, the weight $\gweight_\tau=\vscoef_\tau+\tfrac12\gdiff_\tau^2$, and the source.}
    \label{tab:samplers}
    \small
    \begin{tabular}{p{0.2\linewidth}|p{0.24\linewidth}|p{0.24\linewidth}|p{0.21\linewidth}}
        \toprule
        \textbf{Property} & \textbf{Guided ODE (FM)} & \textbf{SI-SDE} & \textbf{DM-SDE} \\
        \midrule
        Dynamics & deterministic ODE & stochastic SDE & stochastic SDE \\
        Diffusion $\gdiff_\tau$ & $0$ & native $\gamma_\tau$ & lifted $\gdiff_\tau>0$ \\
        Guidance weight $\gweight_\tau$ & $\vscoef_\tau$ & $\vscoef_\tau+\tfrac12\gamma_\tau^2$ & $\vscoef_\tau+\tfrac12\gdiff_\tau^2$ \\
        Source $p_0$ & $\mathcal{N}(0,\bI)$ & $\delta_{\bx_0}$ & $\mathcal{N}(0,\bI)$ \\
        Reweighted init & not needed ($\Phi_0^{\obs}$ const) & vacuous (point mass) & not needed ($\Phi_0^{\obs}$ const) \\
        Velocity/score & from $\fmvel_\theta$ \eqref{eq:fm_score} & from $\drift_\theta$ \eqref{eq:SI_score} & from $\fmvel_\theta$ \eqref{eq:fm_score} \\
        Score recovery & not needed (ODE) & built into drift & required (lift) \\
        \bottomrule
    \end{tabular}
\end{table}


% The construction is exact under a small set of explicit assumptions, collected here with the regime in which each holds; the costliest buy accuracy only where the cheaper ones fail. Table~\ref{tab:approximations} (Appendix~\ref{appendix:bookkeeping}) summarises them by pipeline stage.

% \emph{Modelling.} $\obsoperator$ is linear with additive Gaussian noise (Lemma~\ref{lem:interpolated_likelihood}; exact for linear $\obsoperator$, mild nonlinearity by local linearisation), and the autoregressive factorisation assimilates one observation per step (exact for a single step; otherwise the optimal particle-filter proposal with importance weights and future conditioning omitted).

% \emph{Likelihood surrogate.} The intractable tilt $\Phi_\tau^{\obs}$ is replaced by the observation interpolant (Section~\ref{subsec:obs_interpolation}), with the interpolant likelihood \eqref{eq:interpolant_likelihood_gaussian} taken Gaussian with inflated covariance $\bar\Sigma_\tau=\beta_\tau^2\obscov+\obsmatrix\srccov_\tau\obsmatrix^\top$ \eqref{eq:cov_correction}; covariances are detached when differentiating. This is exact when $p(\bx_1\mid\bx_\tau,\bx_0)$ is Gaussian (Appendix~\ref{appendix:simple_test_case}). The uninflated (DPS) covariance $\beta_\tau^2\obscov$ suffices for dense observations, where the likelihood dominates; for sparse observations the inflation is markedly more accurate, propagating prior-covariance information into the unobserved degrees of freedom (Section~\ref{sec:results}).



% \emph{Schedule and numerics.} Schedule quantities are kept general in $(\alpha_\tau,\beta_\tau,\gamma_\tau)$ (no rectified-flow reduction); the DM-SDE lift uses an endpoint-vanishing $\gdiff_\tau\propto\sqrt{\alpha_\tau\beta_\tau}$ to keep the recovered score finite at $\tau\to1$ (where \eqref{eq:fm_score} diverges); and the correction is clamped to $\tau\in[\Delta\tau,1-\Delta\tau]$ since the coefficients degenerate at the endpoints.

% \subsection{Summary of choices}
% \label{subsec:summary_choices}

% Two orthogonal axes set the configuration: the sampler ($\gdiff_\tau$: ODE / SI-SDE / DM-SDE) and the covariance (Jacobian-free / full / ensemble-shared inflated). Dense observations favour the Jacobian-free covariance, where the samplers nearly coincide; sparse observations favour the ensemble-shared inflated covariance; high-dimensional states are handled through the $N_y\times N_y$ system. The deterministic ODE is cheapest per step, while the SDE samplers give better-calibrated spread. Figure~\ref{fig:decision_guide} condenses these into a decision guide.

% \begin{corollary}[Jacobian-free covariance]\label{cor:cheap_drift}
% Dropping the Jacobian term in $\srccov_\tau$ (i.e.\ $\srccov_\tau\approx\gamma_\tau^2\tau\bI$ for SI and $\srccov_\tau\approx\alpha_\tau^2\bI$ for FM, so $\srccov_\tau\approx\rho_\tau\bI$) replaces the inflated covariance by the state-independent, precomputable
% \begin{align}\label{eq:cov_cheap}
%     \bar{\Sigma}_\tau \approx \beta_\tau^2\,\obscov + \rho_\tau\,\obsmatrix\obsmatrix^\top,
%     \qquad
%     \rho_\tau = \begin{cases} \gamma_\tau^2\tau & \text{(SI)},\\[2pt] \alpha_\tau^2 & \text{(FM)},\end{cases}
% \end{align}
% which depends on $\tau$ only through the scalars $\beta_\tau^2,\rho_\tau$ and requires no network Jacobian. When $\obsmatrix$ and $\obscov$ are time-independent, $\obsmatrix\obsmatrix^\top$ is precomputed once and the per-step solve against $\bar{\Sigma}_\tau$ is cheap.
% \end{corollary}

% Dropping the Jacobian replaces $\srccov_\tau$ by its isotropic part, accurate when $\sigpath_\tau^2\|\nabla_{\bx}^2\log p_\tau\|\ll1$ (near $\tau\to1$, or near-Gaussian priors). For multimodal priors at intermediate $\tau$ the discarded curvature matters and the full inflated covariance may be preferable; we treat this as an empirical trade-off (Section~\ref{sec:results}).

% \begin{remark}[Boundary behaviour]
% At $\tau\to1$ the source scales vanish while $\beta_1=1$, so $\srccov_\tau\to0$ and $\bar{\Sigma}_\tau\to\obscov$: the interpolant likelihood reduces to the raw observation likelihood, consistent with $\interpolantobs_1=\obs$. For $\tau<1$ the inflation enlarges $\bar{\Sigma}_\tau$ along directions visible to $\obsmatrix$, with strength $\rho_\tau/\beta_\tau^2$ relative to $\obscov$.
% \end{remark}

% \subsection{Comparison of posterior samplers}
% \label{subsec:summary_table}\label{subsec:supplying_diffusion}

% The family \eqref{eq:interpolant_posterior_drift} contains three samplers of practical interest, using the \emph{identical} observation-interpolant likelihood with inflated covariance (Section~\ref{subsec:obs_interpolation}), differing only in $\gdiff_\tau$, the weight $\gweight_\tau=\vscoef_\tau+\tfrac12\gdiff_\tau^2$, and the source/anchor.

% Substituting the corresponding diffusion into \eqref{eq:three_drifts} gives each sampler. The \emph{SI-SDE} ($\gdiff_\tau=\gamma_\tau$) reuses the native SDE \eqref{eq:SI_SDE}, whose drift $\drift_\theta=\fmvel_\tau+\tfrac12\gamma_\tau^2\genscore_\tau$ already contains the score, with point-mass source $p_0=\delta_{\bx_0}$ (no reweighting). The \emph{DM-SDE} ($\gdiff_\tau>0$) lifts the deterministic FM ODE \eqref{eq:fm_ode} by recovering the score via \eqref{eq:fm_score}; an endpoint-vanishing $\gdiff_\tau\propto\sqrt{\alpha_\tau\beta_\tau}$ keeps the lifted drift finite as $\tau\to1$ (where \eqref{eq:fm_score} diverges), and the source $\mathcal{N}(0,\bI)$ initialises exactly since $\Phi_0^{\obs}$ is constant (Appendix~\ref{appendix:proof_conditioning}). The \emph{guided ODE} ($\gdiff_\tau=0$) is the probability-flow ODE of the posterior path \eqref{eq:posterior_marginal}, with weight $\vscoef_\tau$; its guidance term coincides with the interpolant-guided flow for linear inverse problems \cite{yan_fig_2024, pokle_training-free_2024}. Being deterministic, all randomness enters through the source -- well posed for FM but not for SI, whose point-mass source would map to a single point.

% Table~\ref{tab:samplers} contrasts the three samplers -- members $\gdiff_\tau=0$ (guided ODE), $\gdiff_\tau=\gamma_\tau$ (SI-SDE), and $\gdiff_\tau>0$ (DM-SDE) of the single family \eqref{eq:unified_posterior_drift} -- differing only in the diffusion, the weight, and the source. The model-specific bookkeeping (source process, how velocity and score are read off the trained network, and the resulting moments and inflated covariance $\bar\Sigma_\tau$) is collected in Table~\ref{tab:si_vs_fm} (Appendix~\ref{appendix:bookkeeping}). The posterior construction -- conditioned path, observation interpolant, and inflated-covariance likelihood score -- is identical across samplers.

% \begin{table}[ht]
%     \centering
%     \caption{The three posterior samplers as members of the unified family \eqref{eq:unified_posterior_drift}. All share the prior drift $\drift^{\gdiff}_\tau=\fmvel_\tau+\tfrac12\gdiff_\tau^2\genscore_\tau$, the likelihood-score correction $\gweight_\tau\,\bar S_\tau$, and the observation-interpolant machinery; they differ only in the diffusion $\gdiff_\tau$, the guidance weight $\gweight_\tau=\vscoef_\tau+\tfrac12\gdiff_\tau^2$, and the source.}
%     \label{tab:samplers}
%     \small
%     \begin{tabular}{p{0.2\linewidth}|p{0.24\linewidth}|p{0.24\linewidth}|p{0.21\linewidth}}
%         \toprule
%         \textbf{Property} & \textbf{Guided ODE (FM)} & \textbf{SI-SDE} & \textbf{DM-SDE} \\
%         \midrule
%         Dynamics & deterministic ODE & stochastic SDE & stochastic SDE \\
%         Diffusion $\gdiff_\tau$ & $0$ & native $\gamma_\tau$ & lifted $\gdiff_\tau>0$ \\
%         Guidance weight $\gweight_\tau$ & $\vscoef_\tau$ & $\vscoef_\tau+\tfrac12\gamma_\tau^2$ & $\vscoef_\tau+\tfrac12\gdiff_\tau^2$ \\
%         Source $p_0$ & $\mathcal{N}(0,\bI)$ & $\delta_{\bx_0}$ & $\mathcal{N}(0,\bI)$ \\
%         Reweighted init & not needed ($\Phi_0^{\obs}$ const) & vacuous (point mass) & not needed ($\Phi_0^{\obs}$ const) \\
%         Velocity/score & from $\fmvel_\theta$ \eqref{eq:fm_score} & from $\drift_\theta$ \eqref{eq:SI_score} & from $\fmvel_\theta$ \eqref{eq:fm_score} \\
%         Score recovery & not needed (ODE) & built into drift & required (lift) \\
%         \bottomrule
%     \end{tabular}
% \end{table}

% \subsection{Approximations and simplifications}
% \label{subsec:approximations}

% The construction is exact under a small set of explicit assumptions, collected here so the scope and the empirical trade-offs of Section~\ref{sec:results} are transparent; each paragraph states the \emph{true quantity}, the \emph{approximation}, and when it is \emph{exact}. Table~\ref{tab:approximations} summarises them, grouped by pipeline stage. Which matter depends on the observation regime, and the costliest approximations buy accuracy only where the cheap ones fail.

% \emph{Modelling.} We take $\obsoperator$ linear with additive Gaussian noise (Lemma~\ref{lem:interpolated_likelihood}; exact for linear $\obsoperator$, mild nonlinearity by local linearisation) and the autoregressive factorisation that assimilates one observation per step (exact for a single step; otherwise the optimal particle-filter proposal with importance weights and future conditioning omitted).

% \emph{Likelihood surrogate.} The intractable tilt $\Phi_\tau^{\obs}$ is replaced by the observation interpolant (Section~\ref{subsec:obs_interpolation}) and the interpolant likelihood \eqref{eq:interpolant_likelihood_gaussian} modelled as a Gaussian with the \emph{inflated} covariance $\bar\Sigma_\tau=\beta_\tau^2\obscov+\obsmatrix\srccov_\tau\obsmatrix^\top$ \eqref{eq:inflated_covariance}; covariances are detached when differentiating. This is exact when $p(\bx_1\mid\bx_\tau,\bx_0)$ is Gaussian (Appendix~\ref{appendix:simple_test_case}, where the inflated covariance is exact). \emph{When it matters:} the inflation is the method's accuracy -- the cheaper uninflated (DPS) covariance $\beta_\tau^2\obscov$ suffices for dense observations, where the likelihood dominates, but for \emph{sparse} observations the inflation is markedly more accurate, as it propagates the prior-covariance information into the unobserved degrees of freedom (Section~\ref{sec:results}).

% \emph{Covariance approximation.} The inflated covariance $\bar\Sigma_\tau$ carries a network-Jacobian term in $\srccov_\tau$. The \emph{Jacobian-free} covariance replaces $\srccov_\tau$ by $\rho_\tau\bI$ (Corollary~\ref{cor:cheap_drift}), giving the precomputable $\bar\Sigma_\tau\approx\beta_\tau^2\obscov+\rho_\tau\obsmatrix\obsmatrix^\top$, accurate when $\sigpath_\tau^2\|\nabla_{\bx}^2\log p_\tau\|\ll1$ -- the method of choice for dense observations and large $N_y$. Retaining the full $\srccov_\tau$ needs $N_y$ Jacobian--vector products \emph{per ensemble member} per step, $\mathcal{O}(E\,N_y)$ and prohibitive at field scale; the \emph{ensemble-shared} approximation evaluates the source-covariance Jacobian once at the ensemble mean, factorising a single $N_y\times N_y$ system and reducing this to $\mathcal{O}(N_y)$, exact as the ensemble tightens. This EnKF-analogue step recovers accuracy in the sparse regime where the Jacobian-free covariance degrades (Section~\ref{sec:results}).

% \emph{Schedule and numerics.} All schedule-derived quantities are kept general in $(\alpha_\tau,\beta_\tau,\gamma_\tau)$ (no rectified-flow reduction); the DM-SDE lift uses an endpoint-vanishing $\gdiff_\tau\propto\sqrt{\alpha_\tau\beta_\tau}$ to keep the recovered score finite at $\tau\to1$; and the correction is clamped to $\tau\in[\Delta\tau,1-\Delta\tau]$ since the coefficients degenerate at the endpoints. Table~\ref{tab:approximations} (Appendix~\ref{appendix:bookkeeping}) enumerates all of these with the regime in which each is exact.

% \subsection{Summary of choices}
% \label{subsec:summary_choices}

% The framework has two orthogonal axes -- the \emph{sampler} ($\gdiff_\tau$: ODE / SI-SDE / DM-SDE) and the \emph{covariance} (Jacobian-free / full / ensemble-shared inflated) -- chosen by problem geometry: dense observations favour the cheap Jacobian-free covariance (the samplers nearly coincide), sparse observations the ensemble-shared inflated covariance; high-dimensional states are handled through the $N_y\times N_y$ system; and the deterministic ODE is cheapest per step while SDE samplers give better-calibrated spread. Figure~\ref{fig:decision_guide} condenses these into a single decision guide.

% \input{figures/decision_guide}
